\section{Per-airport runway thresholds and ARP coordinates}
\label{sec:appendix-airports}

The per-airport runway thresholds and ARP elevations used by the SDF
builder of \cref{sec:method-sdf} are taken verbatim from the
public-domain FAA NASR Airport Diagrams effective in August 2024.
Reproducing the eight-row table per airport here would duplicate the
project's \texttt{configs/airports/\{KLAX,KSFO\}.yaml} files; we point
the reader to those YAML files instead.

% Per-airport runway tables omitted from this preprint draft. They are
% available verbatim in the project repository at
% configs/airports/{KLAX,KSFO}.yaml and consist of 8 rows per airport
% (threshold lat/lon, end lat/lon, length, width, true bearing,
% code letter/number, precision flag).

\section{Annex~14 OLS parameter table}
\label{sec:appendix-ols}

The complete OLS parameter set used in \cref{sec:method-sdf} appears in the
body as \cref{tab:ols-params}. Per-airport overrides are supported in
\texttt{configs/annex14/} but are unused in this paper -- both LAX and SFO
are evaluated with the identical Code~4 precision approach parameter set.

\section{Counterfactual sampling details}
\label{sec:appendix-cf}

For each (date, 15-min slice) pair we draw $n$ candidate eVTOL segments
uniformly from the per-vertiport OFV $\Vofv$ of the source vertiport
$v_{\mathrm{src}}$ to a uniform random destination inside the SDF box
intersected with $\Astatic$. The segment kinematic profile is a piecewise
linear trajectory: climb at $7$~m/s to cruise altitude (a uniform random
choice in $[300, 1000]$~m AGL), level cruise at $67$~m/s (cruise speed
from \texttt{configs/scenarios/cost\_weights.yaml}), descent at $5$~m/s to
the destination OFV. The segment is discretised at $1$~s and matched
against every contemporaneous ADS-B state vector within
$\pm 60$~s for the five conflict criteria listed in
\cref{sec:method-risk}.

For LAX and SFO we generate $50\,000$~/~$40\,000$~/~$40\,000$~/~$40\,000$~/~
$40\,000$ segments per Friday respectively, for a total of $210\,000$ LAX
and $210\,000$ SFO segments. We hold out 2024-08-30 as the temporal test
fold.

\section{Computational footprint}
\label{sec:appendix-compute}

The entire pipeline was executed on a single Featurize RTX~4090 49\,GB
instance (Featurize, Hangzhou). One-time stage costs:

\begin{description}
    \item[OLS SDF build (one airport)] $\sim 45$~s on a single CPU thread.
    \item[ADS-B clean + classify (one airport-day)] $\sim 90$~s on a
        single CPU thread.
    \item[Dynamic-envelope build (one airport-day)] $\sim 8$~s.
    \item[Counterfactual sampling (50k segs)] $\sim 4$~min on a single
        CPU thread; embarrassingly parallel across days.
    \item[Feature extraction (220k segs across 5 days)] $\sim 40$~min
        on a single CPU thread, dominated by an $O(N_{\mathrm{seg}}
        \times N_{\mathrm{ADS\text{-}B}})$ density-box scan; a
        vectorised replacement is straightforward and is part of the
        future work in \cref{sec:discussion}.
    \item[Risk-field training (XGB on cached features)] $\sim 7$~s once
        features are cached.
    \item[Corridor planning (180 corridors per airport-baseline at
        $300\times90$~m planning resolution)] $\sim 8$~min on 5 parallel
        CPU threads.
\end{description}

The dominant single operation is feature extraction at $40$~min per
airport. Total wall-clock time from raw ADS-B to final figures is under
$3$~hours per airport.
